\section{Introduction}

Despite decades of work on connecting arguments across the web---from the Semantic Web \citep{Berners-Lee:2001} to the Argument Web \citep{Bex:2013} and AIF \citep{Chesnevar:2007:AIF}---argumentation software has not reached the mainstream. Most argument enthusiasts still turn to social media platforms like Reddit, Twitter, and Substack, whose success stems from open, flexible participation with feedback \citep{Leite2011}, delivered through simple text-based interfaces where arguments flow naturally in written form. Meanwhile, the argument mining community has produced increasingly capable models for extracting argumentative structure from text, and there is a growing opportunity to translate these capabilities into applications.

We explore whether commercially available LLMs, with only few-shot prompting, can perform the end-to-end argument mining needed to power such a platform: extracting argumentative discourse units (ADUs), canonicalising and deduplicating them across discussions, and using the resulting structure to rank contributions by argument quality rather than just social signals. We evaluate each pipeline stage on established benchmarks---Persuasive Essays \citep{stab-gurevych-2014-annotating} for extraction, MultiPIT \citep{dou-etal-2022-improving} for deduplication---and show through an end-to-end evaluation on ChangeMyView \citep{Kolyada2020} that even the noisy results LLMs produce have enough useful structure to significantly improve reply ranking.

A platform that ranks by argument quality rather than popularity is particularly relevant as social media becomes increasingly populated by automated accounts, with estimates of 9--15\% prevalence among active users \citep{varol2017online} and substantially higher rates during crisis events \citep{elayan2024digital}. Conventional platforms prioritize engagement over quality, allowing bots to steer discussion through volume alone \citep{farcas2025manufacturing}. Yet, multi-agent LLM debate has shown that structured interaction between machine agents can surface implicit premises and improve reasoning quality \citep{ku2025multiagent}. A ranking algorithm grounded in argument structure---weighting \emph{what} is argued over \emph{how much} engagement it receives---is architecturally resilient to such manipulation: low-quality arguments rank low regardless of how many accounts produce them. This property makes argument-aware ranking a prerequisite for platforms that aim to compile reliable knowledge from both human and machine reasoning.

\paragraph{Related work.} ChangeMyView has been widely studied as a natural laboratory for persuasion. \citet{tan-etal-2016-winning} established the use of delta awards as a persuasion signal and showed that interaction dynamics and linguistic features predict persuasion success. \citet{wei-etal-2016-post} ranked CMV comments using argumentation features---including counts of argumentative sentences---and found them more predictive than surface text features. \citet{dutta-etal-2019-changing} combined persuasion modelling with semi-supervised argument extraction from CMV threads. More broadly, \citet{lawrence-reed-2017-complex} showed that structural metrics over large-scale argument graphs from online political debates yield informative features. LLMs have been increasingly applied to argument mining \citep[for a recent survey, see][]{li-etal-2025-llm-am-survey}, though primarily in classification settings. \citet{chen-etal-2024-exploring-potential} systematically evaluated LLMs across six argumentation tasks in zero-shot and few-shot settings, finding promising but sub-SOTA performance; \citet{alzubaer-etal-2023-legal-am} reported similar findings for argument component classification in legal text, where domain-specific models still outperform GPT-4. Other work has shown that zero-shot GPT models can detect fallacies above baseline \citep{ruiz-dolz-lawrence-2023-gpt}, Chain-of-Thought prompting improves on structured text \citep{bytyqi-hautli-janisz-2025-reasoning}, and few-shot priming aids stance classification \citep{ajjour-wachsmuth-2025-exploring}. LLM-based parsers have also been applied to dialogical argument mining \citep{saha-srihari-2024-turiya} and argument detection under resource constraints \citep{kashefi-etal-2023-argument}. Beyond classification, LLMs have been used for interactive argument graph construction \citep{lenz-bergmann-2024-arguemapper}, multi-task argumentation via fine-tuning \citep{savingy-yun-2025-amelia}, and combined with formal argumentation semantics for explainable debate analysis \citep{alfano-etal-2025-fakb}. Mapping arguments to canonical forms is a recurring challenge addressed by Key Point Analysis \citep{bar-haim-etal-2020-arguments}, hybrid embedding-plus-human pipelines \citep{vandermeer-etal-2024-hyena}, and embedding clustering for web-scale deduplication \citep{abbas2023semdedup}. We take a complementary approach, deploying LLMs as the engine for a complete pipeline---from extraction through deduplication to ranking---and evaluating whether the resulting structure improves a downstream application.
